% GENERATED FILE. Edit canonical lessons, then run npm run generate:books.
% canonical-source-hash: fnv1a64:7a760236990190c6
% canonical-lessons: ZH-C229-chongwu, ZH-C229-yang, ZH-C229-liugou, ZH-C229-zhua, ZH-C229-mo, ZH-R229-review-224-229

\chapter{Keeping a Pet}
\label{ch:zh-keeping-a-pet}
\hlchaptermodality{229}

\begin{chapteropening}
\textbf{By the end of this chapter:} \emph{I can talk about pets: keeping one, walking the dog, being scratched and stroking them.}

Complete the last lesson of chapter 229: I can talk about pets: keeping one, walking the dog, being scratched and stroking them.
\end{chapteropening}

\section[chǒngwù]{chǒngwù — pet}

\label{lesson:ZH-C229-chongwu}



\begin{quote}
Before the new one: say the Mandarin for serious (of illness), then the Mandarin for to nod.
\end{quote}

\subsection*{You'll want to know: chǒngwù}
\emph{chǒngwù} — \textquotedblleft{}pet\textquotedblright{}.

\begin{quote}
\emph{nǐ jiā yǒu chǒngwù ma} — Have you got a pet at home?
\end{quote}

Its characters come later; for now it is a word for the ear.

\begin{cousinweb}
\emph{chǒngwù} is built from \emph{chǒng} 'to dote on' and \emph{wù} 'thing', both inherited from Middle Chinese, \emph{trhjowngX} and \emph{mjut} (Guangyun, in Baxter's transcription), which glosses the first 'to love and favour'.
\end{cousinweb}

\begin{grammarlens}[title={What you've built}]
One more word for saying what I do.
\end{grammarlens}

\subsection*{Your turn}
\begin{itemize}
\raggedright
  \item \emph{Say it:} \emph{chǒngwù}
  \item \emph{Say it:} \emph{chǒngwù}, once more
  \item \emph{Say it:} \emph{nǐ jiā yǒu chǒngwù ma} — Have you got a pet at home?
  \item \emph{Recall:} say the Mandarin for serious (of illness), then the Mandarin for to nod, then say \emph{chǒngwù} again
\end{itemize}

\begin{tcolorbox}[breakable,colback=teal!4,colframe=teal!35!black,title={Before you move on}]
What does chǒngwù mean? (\textquotedblleft{}Pet\textquotedblright{}.) Say it once more.
\end{tcolorbox}
\section[yǎng]{yǎng — to keep an animal; to raise}

\label{lesson:ZH-C229-yang}



\begin{quote}
Before the new one: say the Mandarin for to nod, then the Mandarin for a pet.
\end{quote}

\subsection*{You'll want to know: yǎng}
\emph{yǎng} — \textquotedblleft{}to keep an animal; to raise\textquotedblright{}.

\begin{quote}
\emph{wǒ xiǎng yǎng yì zhī māo} — I'd like to have a cat.
\end{quote}

Its characters come later; for now it is a word for the ear.

\begin{cousinweb}
\emph{yǎng} is inherited from Middle Chinese \emph{yangX} (Guangyun, in Baxter's transcription), glossed 'to rear' and 'to delight'.
\end{cousinweb}

\begin{grammarlens}[title={What you've built}]
One more word for saying what I do.
\end{grammarlens}

\subsection*{Your turn}
\begin{itemize}
\raggedright
  \item \emph{Say it:} \emph{yǎng}
  \item \emph{Say it:} \emph{yǎng}, once more
  \item \emph{Say it:} \emph{wǒ xiǎng yǎng yì zhī māo} — I'd like to have a cat.
  \item \emph{Recall:} say the Mandarin for to nod, then the Mandarin for a pet, then say \emph{yǎng} again
\end{itemize}

\begin{tcolorbox}[breakable,colback=teal!4,colframe=teal!35!black,title={Before you move on}]
What does yǎng mean? (\textquotedblleft{}To keep an animal; to raise\textquotedblright{}.) Say it once more.
\end{tcolorbox}
\section[liùgǒu]{liùgǒu — to walk the dog}

\label{lesson:ZH-C229-liugou}



\begin{quote}
Before the new one: say the Mandarin for a pet, then the Mandarin for to keep an animal.
\end{quote}

\subsection*{You'll want to know: liùgǒu}
\emph{liùgǒu} — \textquotedblleft{}to walk the dog\textquotedblright{}.

\begin{quote}
\emph{wǒ měitiān wǎnshang liùgǒu} — I walk the dog every evening.
\end{quote}

Its characters come later; for now it is a word for the ear.

\begin{cousinweb}
\emph{liùgǒu} is built from \emph{liù} 'to stroll' and \emph{gǒu} 'dog'. The second is inherited from Middle Chinese \emph{kuwX} (Guangyun, in Baxter's transcription), glossed 'dog'; the dictionary glosses the first, \emph{ljuw} there, 'to linger', and the falling tone of liù is not the regular outcome of that form.
\end{cousinweb}

\begin{grammarlens}[title={What you've built}]
One more word for saying what I do.
\end{grammarlens}

\subsection*{Your turn}
\begin{itemize}
\raggedright
  \item \emph{Say it:} \emph{liùgǒu}
  \item \emph{Say it:} \emph{liùgǒu}, once more
  \item \emph{Say it:} \emph{wǒ měitiān wǎnshang liùgǒu} — I walk the dog every evening.
  \item \emph{Recall:} say the Mandarin for a pet, then the Mandarin for to keep an animal, then say \emph{liùgǒu} again
\end{itemize}

\begin{tcolorbox}[breakable,colback=teal!4,colframe=teal!35!black,title={Before you move on}]
What does liùgǒu mean? (\textquotedblleft{}To walk the dog\textquotedblright{}.) Say it once more.
\end{tcolorbox}
\section[zhuā]{zhuā — to scratch; to grab}

\label{lesson:ZH-C229-zhua}



\begin{quote}
Before the new one: say the Mandarin for to keep an animal, then the Mandarin for to walk the dog.
\end{quote}

\subsection*{You'll want to know: zhuā}
\emph{zhuā} — \textquotedblleft{}to scratch; to grab\textquotedblright{}.

\begin{quote}
\emph{bié pà wǒ de māo bù zhuā rén} — Don't worry, my cat doesn't scratch.
\end{quote}

Its characters come later; for now it is a word for the ear.

\begin{cousinweb}
\emph{zhuā} is inherited from Middle Chinese \emph{tsraew} (Guangyun, in Baxter's transcription), glossed 'to scratch and pinch'; the vowel of zhuā is not the regular outcome of that form.
\end{cousinweb}

\begin{grammarlens}[title={What you've built}]
One more word for saying what I do.
\end{grammarlens}

\subsection*{Your turn}
\begin{itemize}
\raggedright
  \item \emph{Say it:} \emph{zhuā}
  \item \emph{Say it:} \emph{zhuā}, once more
  \item \emph{Say it:} \emph{bié pà wǒ de māo bù zhuā rén} — Don't worry, my cat doesn't scratch.
  \item \emph{Recall:} say the Mandarin for to keep an animal, then the Mandarin for to walk the dog, then say \emph{zhuā} again
\end{itemize}

\begin{tcolorbox}[breakable,colback=teal!4,colframe=teal!35!black,title={Before you move on}]
What does zhuā mean? (\textquotedblleft{}To scratch; to grab\textquotedblright{}.) Say it once more.
\end{tcolorbox}
\section[mō]{mō — to stroke; to touch}

\label{lesson:ZH-C229-mo}



\begin{quote}
Before the new one: say the Mandarin for to walk the dog, then the Mandarin for to scratch.
\end{quote}

\subsection*{You'll want to know: mō}
\emph{mō} — \textquotedblleft{}to stroke; to touch\textquotedblright{}.

\begin{quote}
\emph{wǒ kěyǐ mō nǐ de gǒu ma} — Can I stroke your dog?
\end{quote}

Its characters come later; for now it is a word for the ear.

\begin{cousinweb}
\emph{mō} is inherited from Middle Chinese \emph{mu} (Guangyun, in Baxter's transcription), glossed 'to feel with the hand'. The vowel and the level tone of mō are not the regular outcome of that form.
\end{cousinweb}

\begin{grammarlens}[title={What you've built}]
One more word for saying what I do.
\end{grammarlens}

\subsection*{Your turn}
\begin{itemize}
\raggedright
  \item \emph{Say it:} \emph{mō}
  \item \emph{Say it:} \emph{mō}, once more
  \item \emph{Say it:} \emph{wǒ kěyǐ mō nǐ de gǒu ma} — Can I stroke your dog?
  \item \emph{Recall:} say the Mandarin for to walk the dog, then the Mandarin for to scratch, then say \emph{mō} again
\end{itemize}

\begin{tcolorbox}[breakable,colback=teal!4,colframe=teal!35!black,title={Before you move on}]
What does mō mean? (\textquotedblleft{}To stroke; to touch\textquotedblright{}.) Say it once more.
\end{tcolorbox}
\section[(dialogue)]{Looking back: from Spring Flowers to Keeping a Pet}

\label{lesson:ZH-R229-review-224-229}



\begin{quote}
Say the Mandarin for to bloom and a florist's.
\end{quote}

\subsection*{Your turn}
\begin{itemize}
\raggedright
  \item \emph{Say it:} to bloom, a florist's, a garden, a butterfly, a bee
  \item \emph{Say it:} the countryside, a farmer, to work the land, to pick fruit, to camp out
  \item \emph{Say it:} a lawn, a lake, a forest, grassland, scenery
  \item \emph{Say it:} the seaside, a boat, to get seasick, a sandy beach, to run into someone
  \item \emph{Say it:} a headache, to take medicine, to have an injection, serious (of illness), to nod
  \item \emph{Say it:} a pet, to keep an animal, to walk the dog, to scratch, to stroke
  \item \emph{Read it:} \zh{花} \zh{草} \zh{头} \zh{疼} — each one aloud, then a word that holds it
\end{itemize}

\begin{tcolorbox}[breakable,colback=teal!4,colframe=teal!35!black,title={Before you move on}]
To bloom? (\textbf{\zh{开花}}, \emph{kāihuā}.) A florist's? (\textbf{\zh{花店}}, \emph{huādiàn}.) A garden? (\emph{huāyuán}.) A butterfly? (\emph{húdié}.) A bee? (\emph{mìfēng}.) The countryside? (\emph{nóngcūn}.) A farmer? (\emph{nóngmín}.) To work the land? (\emph{zhòngdì}.) To pick fruit? (\emph{zhāi}.) To camp out? (\emph{lùyíng}.) A lawn? (\textbf{\zh{草地}}, \emph{cǎodì}.) A lake? (\emph{hú}.) A forest? (\emph{sēnlín}.) Grassland? (\emph{cǎoyuán}.) Scenery? (\emph{fēngjǐng}.) The seaside? (\emph{hǎibiān}.) A boat? (\emph{chuán}.) To get seasick? (\emph{yùnchuán}.) A sandy beach? (\emph{shātān}.) To run into someone? (\emph{yùdào}.) A headache? (\textbf{\zh{头疼}}, \emph{tóuténg}.) To take medicine? (\emph{chīyào}.) To have an injection? (\emph{dǎzhēn}.) Serious (of illness)? (\emph{yánzhòng}.) To nod? (\textbf{\zh{点头}}, \emph{diǎntóu}.) A pet? (\emph{chǒngwù}.) To keep an animal? (\emph{yǎng}.) To walk the dog? (\emph{liùgǒu}.) To scratch? (\emph{zhuā}.) To stroke? (\emph{mō}.)
\end{tcolorbox}
